\exercisesection{\S~4}

\begin{enumerate}
\def\labelenumi{\arabic{enumi})}
\tightlist
\item
  Take G to be one of the groups \(\mathbf{SU}(n,\mathbf{C})\) or \(\mathbf{U}(n,\mathbf{H})\) ; identify \(g_{C}\) with \(\mathfrak{sl}(n,\mathbf{C})\) or \(\mathfrak{sp}(2n,\mathbf{C})\) , respectively, cf.~§3, no. 4 and Exerc. 3. We use the notations of Chap. VIII, §13, nos. 1 and 3, with k=C.
\end{enumerate}

\begin{enumerate}
\def\labelenumi{\alph{enumi})}
\item
  Show that the subgroup T of G consisting of the diagonal matrices with complex coefficients is a maximal torus, and that \(\mathrm{L}(\mathrm{T})_{(\mathbf{C})} = \mathfrak{h}\) .
\item
  Identify X(T) with a subgroup of \(\mathfrak{h}^*\) by means of the homomorphism \(\delta\) . Show that the linear forms \(\varepsilon_i\) and the dominant weights \(\varpi_j\) belong to X(T). If \(t = \mathrm{diag}(t_1, \ldots, t_n) \in \mathrm{T}\) , then \(\varepsilon_i(t) = t_i\) and \(\varpi_j(t) = t_1 \ldots t_j\) for \(1 \leq i, j \leq n\) .
\item
  Deduce from b) another proof of the fact that the groups \(\mathbf{SU}(n,\mathbf{C})\) and \(\mathbf{U}(n,\mathbf{H})\) are simply-connected (cf.~§3, Exerc. 4 and 6).
\end{enumerate}

\begin{enumerate}
\def\labelenumi{\arabic{enumi})}
\setcounter{enumi}{1}
\tightlist
\item
  Take \(\mathrm{G} = \mathbf{SO}(n, \mathbf{R})\) , with \(n \geq 3\) ; put \(n = 2l + 1\) if n is odd, and n = 2l if n is even. The algebra \(g_{C}\) can be identified with \(\mathfrak{o}(n, \mathbf{C})\) ; we use the notations of Chap. VIII, §13, nos. 2 and 4. Denote by \((f_i)_{1 \leq i \leq n}\) the canonical basis of \(R^n\) . Put \(e_j = \frac{1}{\sqrt{2}}(f_{2j-1} + if_{2j})\) and \(e_{-j} = \frac{1}{\sqrt{2}}(f_{2j-1} - if_{2j})\) for \(1 \leq j \leq l\) , and \(e_0 = i\sqrt{2}f_{2l+1}\) if n is odd; choose the Witt basis of \(C^n\) given by \(e_1, \ldots, e_l, e_{-l}, \ldots, e_{-1}\) if n is even (resp. \(e_1, \ldots, e_l, e_0, e_{-l}, \ldots, e_{-1}\) if n is odd).
\end{enumerate}

\begin{enumerate}
\def\labelenumi{\alph{enumi})}
\tightlist
\item
  Let \(H_{i}\) be the subspace of \(R^{n}\) generated by \(f_{2i-1}\) and \(f_{2i}\) ( \(1 \leq i \leq l\) ); show that the subgroup of G consisting of the elements g such that \(g(\mathrm{H}_{i}) \subset \mathrm{H}_{i}\) and \(\det(g|\mathrm{H}_{i}) = 1\) for \(1 \leq i \leq l\) is a maximal torus T of G, and that \(\mathrm{L}(\mathrm{T})_{(\mathbf{C})} = \mathfrak{h}\) .
\item
  Identify X(T) with a subgroup of \(h^{*}\) by means of \(\delta\) . Show that the linear forms \(\varepsilon_{i}\) belong to X(T); if \(t \in T\) and if the restriction of t to \(H_{j}\) is a rotation through an angle \(\theta_{j}\) , then \(\varepsilon_{j}(t) = e^{i\theta_{j}}\) . The weights \(\varpi_{1}, \ldots, \varpi_{l-2}\) ; \(2\varpi_{l-1}, 2\varpi_{l}, \varpi_{l-1} \pm \varpi_{l}\) belong to X(T). If n is odd, \(\varpi_{l-1}\) belongs to X(T).
\item
  Let \(\tilde{\mathbf{G}} = \mathbf{Spin}(n, \mathbf{R})\) and let \(\varphi : \tilde{G} \to G\) be the canonical covering. For \(\boldsymbol{\theta} = (\theta_{1}, \ldots, \theta_{l}) \in \mathbf{R}^{l}\) , put \(t(\boldsymbol{\theta}) = \prod_{i=1}^{l} (\cos\theta_{i} - f_{2i-1}f_{2i}\sin\theta_{i}) \in \tilde{\mathbf{G}}\) . Show that the set of the \(t(\boldsymbol{\theta})\) for \(\theta \in R^{l}\) is a maximal torus \(\tilde{T}\) of \(\tilde{G}\) , such that \(\varphi(\tilde{T}) = T\) . If X( \(\tilde{T}\) ) is identified with a subgroup of \(h^{*}\) , we have \(\varepsilon_{j}(t(\boldsymbol{\theta})) = e^{2i\theta_{j}}\) .
\item
  Show that the weights \(\varpi_{l-1}\) and \(\varpi_{l}\) belong to X( \(\tilde{T}\) ); deduce that \(\mathbf{Spin}(n,\mathbf{R})\) is simply-connected (cf.~§3, Exerc. 5).
\end{enumerate}

\begin{enumerate}
\def\labelenumi{\arabic{enumi})}
\setcounter{enumi}{2}
\tightlist
\item
  \begin{enumerate}
  \def\labelenumii{\alph{enumii})}
  \tightlist
  \item
    Show that the automorphism \(\sigma: \mathrm{A} \mapsto \overline{\mathrm{A}}\) of \(\mathbf{SU}(n, \mathbf{C})\) is not inner for \(n \geq 3\) . Every non-inner automorphism of \(\mathbf{SU}(n, \mathbf{C})\) is of the form \((\operatorname{Int} g) \circ \sigma\) , \(g \in \mathbf{SU}(n, \mathbf{C})\) .
  \end{enumerate}
\end{enumerate}

\begin{enumerate}
\def\labelenumi{\alph{enumi})}
\setcounter{enumi}{1}
\tightlist
\item
  Show that for every almost simple connected compact group G of type \(A_{n}\) ( \(n \geq 2\) ), the group \(\operatorname{Aut}(\mathrm{G})/\operatorname{Int}(\mathrm{G})\) is cyclic of order 2 (cf.~§3, Exerc. 4).
\end{enumerate}

\begin{enumerate}
\def\labelenumi{\arabic{enumi})}
\setcounter{enumi}{3}
\tightlist
\item
  Let \(n\) be an integer \(\geq 2\) .
\end{enumerate}

\begin{enumerate}
\def\labelenumi{\alph{enumi})}
\item
  Let \(g \in \mathbf{O}(2n, \mathbf{R})\) with \(\det g = -1\) ; show that the automorphism \(\operatorname{Int} g\) of \(\mathbf{O}(2n, \mathbf{R})\) induces an automorphism of \(\mathbf{SO}(2n, \mathbf{R})\) that is not inner.
\item
  For \(n \geq 2\) , the group \(\operatorname{Aut}(\mathbf{SO}(2n, \mathbf{R}))\) is equal to \(\operatorname{Int}(\mathbf{O}(2n, \mathbf{R}))\) (which is isomorphic to \(\mathbf{O}(2n, \mathbf{R}) / \{\pm I_{2n}\}\) ).
\item
  Establish an analogous result for \(\mathbf{Spin}(2n,\mathbf{R})\) and \(\mathbf{SO}(2n,\mathbf{R})/\{\pm I_{2n}\}\) . If n is even and \(\neq 2\) , every automorphism of the group \(\mathbf{Spin}(2n,\mathbf{R})/\{1,\varepsilon\}\) (§3, Exerc. 5) is inner.
\end{enumerate}

\begin{enumerate}
\def\labelenumi{\arabic{enumi})}
\setcounter{enumi}{4}
\tightlist
\item
  Let \(\mathbf{R}\) be a reduced and irreducible root system.
\end{enumerate}

\begin{enumerate}
\def\labelenumi{\alph{enumi})}
\item
  Show that the set of roots in R of greatest length is a symmetric closed subset of R, stable under W(R).
\item
  Show that every non-empty subset P of R that is closed, symmetric and stable under W(R) is equal to R or to the set of roots of greatest length in R.
\item
  Assume that \(P \neq R\) . Show that the root system P is of type \(D_{l}\) (resp. \((\mathrm{A}_{1})^{l}, \mathrm{D}_{4}, \mathrm{A}_{2}\) ) if R is of type \(B_{l}\) (resp. \(C_{l}, F_{4}, G_{2}\) ).
\end{enumerate}

\begin{enumerate}
\def\labelenumi{\arabic{enumi})}
\setcounter{enumi}{5}
\tightlist
\item
  Let \(\mathrm{H}\) be a closed subgroup of \(\mathrm{G}\) containing \(\mathrm{N}_{\mathrm{G}}(\mathrm{T})\) .
\end{enumerate}

\begin{enumerate}
\def\labelenumi{\alph{enumi})}
\item
  Show that H is equal to its own normalizer in G, and that \(H/H_{0}\) is isomorphic to \(W/W_{H_{0}}(T)\) .
\item
  Show that \(\mathrm{R}(\mathrm{H}_{0},\mathrm{T})\) is stable under W. Conversely, if K is a connected closed subgroup containing T such that \(\mathrm{R}(\mathrm{K},\mathrm{T})\) is stable under W, the normalizer of K contains that of T.
\end{enumerate}

\(c)\) Assume that \(\mathbf{G}\) is almost simple. Prove that \(\mathbf{H}\) is equal to \(\mathrm{N_G(T)}\) , or to \(\mathbf{G}\) , or that we are (up to isomorphism) in one of the following situations:

\(\alpha\) (resp. \(\alpha'\) ) \(\mathrm{G} = \mathbf{Spin}(2l + 1, \mathbf{R})\) (resp. \(\mathrm{G} = \mathbf{SO}(2l + 1, \mathbf{R})\) ) and \(\mathrm{H}_0\) is the fixer of a non-zero vector of \(\mathbf{R}^{2l + 1}\) , isomorphic to \(\mathbf{Spin}(2l, \mathbf{R})\) (resp. \(\mathbf{SO}(2l, \mathbf{R})\) );

\(\beta)\) \(\mathrm{G} = \mathrm{U}(l,\mathbf{H})\) and \(\mathrm{H}_0\) is the subgroup D consisting of the diagonal matrices; \(\beta^{\prime})\) \(\mathrm{G} = \mathrm{U}(l,\mathbf{H}) / \{\pm 1\}\) and \(\mathrm{H}_0 = \mathrm{D} / \{\pm 1\}\) ;

\(\gamma)\) G is of type \(\mathrm{F}_4\) and \(\mathrm{H}_0\) is isomorphic to Spin(8,R);

\(\delta)\) G is of type \(\mathrm{G}_2\) and \(\mathrm{H}_0\) is the subgroup (isomorphic to \(\mathbf{SU}(3,\mathbf{C})\) ) defined in Exerc. 7 of §3.

\begin{enumerate}
\def\labelenumi{\arabic{enumi})}
\setcounter{enumi}{6}
\tightlist
\item
  Let \(\tau : \mathbf{G} \to \mathbf{GL}(\mathbf{V})\) be a continuous representation of \(\mathbf{G}\) on a finite dimensional real vector space. Assume that for all \(\lambda \in \mathrm{X}(\mathrm{T})\) , we have \(\dim_{\mathbf{C}} \hat{V}_{\lambda} \leq 1\) .
\end{enumerate}

\begin{enumerate}
\def\labelenumi{\alph{enumi})}
\item
  Show that the representation \(\tau\) is the direct sum of a finite family \((\tau_{i})_{1\leq i\leq s}\) of irreducible representations, mutually non-isomorphic, whose commutants \(K_{i}\ (1\leq i\leq s)\) are isomorphic to R or C.
\item
  There exists a C-vector space structure on V for which the operations of \(\tau(\mathrm{G})\) are C-linear if and only if \(K_{1},\ldots,K_{s}\) are isomorphic to C; there are then \(2^{s}\) structures of this type.
\end{enumerate}

\begin{enumerate}
\def\labelenumi{\arabic{enumi})}
\setcounter{enumi}{7}
\tightlist
\item
  Let H be a connected closed subgroup of G of maximum rank, h its Lie algebra, X the manifold G/H and V the tangent space of X at the point corresponding to the class of H; identify V with the quotient g/h.
\end{enumerate}

\begin{enumerate}
\def\labelenumi{\alph{enumi})}
\tightlist
\item
  If j is an almost complex structure (Differentiable and Analytic Manifolds, Results, 8.8.3) on X, denote by \(V''(j)\) the subspace of the complex vector space \(C \otimes V\) consisting of the elements u such that \(j(u) = -iu\) , and by \(\mathfrak{q}(j)\) the subspace of \(g_{C}\) that is the inverse image of \(V''(j)\) , so that the canonical map \(V \to \mathfrak{g}_{\mathbf{C}} / \mathfrak{q}(j)\) is a C-linear isomorphism when V is given the C-vector space structure induced by j. Prove that the map \(j \mapsto \mathfrak{q}(j)\) is a bijection from the set of almost complex structures on X invariant under G to the set of complex subspaces p of \(g_{C}\) satisfying the following conditions:
\end{enumerate}

\[
\mathfrak {p} + \bar {\mathfrak {p}} = \mathfrak {g} _ {\mathbf {C}}\tag{1}
\]

\[
\mathfrak {p} \cap \bar {\mathfrak {p}} = \mathfrak {h} _ {\mathbf {C}}\tag{2}
\]

\[
[ \mathfrak {h}, \mathfrak {p} ] \subset \mathfrak {p}.\tag{3}
\]

\begin{enumerate}
\def\labelenumi{\alph{enumi})}
\setcounter{enumi}{1}
\item
  There exists such a structure on X if and only if the commuting fields of the irreducible subrepresentations of the adjoint representation of H on V are all isomorphic to C; in that case, there are \(2^{s}\) such structures, where s is the number of irreducible subrepresentations of V (use Exerc. 7).
\item
  Let j be an almost complex structure on X, invariant under G, and \(\mathfrak{p} = \mathfrak{q}(j)\) the associated subspace. Then, j is integrable (that is, associated to a complex-analytic manifold structure on X, cf.~Differentiable and Analytic Manifolds, Results, 8.8.5 to 8.8.8) if and only if p satisfies the condition
\end{enumerate}

\[
[ \mathfrak {p}, \mathfrak {p} ] \subset \mathfrak {p}.\tag{4}
\]

\begin{enumerate}
\def\labelenumi{\alph{enumi})}
\setcounter{enumi}{3}
\item
  There exists a complex structure (that is, an integrable almost complex structure) on X invariant under G if and only if H is the centralizer of a torus in G (show that conditions (1) to (4) above imply that p is a parabolic subalgebra of \(g_{C}\) (Chap. VIII, §3, no. 5)); in that case, these complex structures correspond bijectively to the parabolic subalgebras p of \(g_{C}\) that are the direct sum of \(h_{C}\) and their unipotent radical (cf.~Chap. VIII, §3, no. 4).
\item
  Prove that there exist exactly \(\operatorname{Card}(W)\) complex structures on G/T invariant under G; if \(\sigma\) and \(\sigma'\) are two such structures, there exists a unique element \(w \in W\) such that the canonical operation of w on G/T (by inner automorphisms) transforms \(\sigma\) into \(\sigma'\) . If w is a non-identity element of W, the operation of w on G/T is not C-analytic for any complex structure on G/T invariant under G.
\item
  Determine the complex structures on \(S^{2}\) invariant under \(\mathbf{SO}(3)\) .
\end{enumerate}

\(g)^{*}\) With the notations of Exerc. 8 \(d\) ), let \(\mathrm{G}_c\) be the complexification of \(\mathrm{G}\) and \(\mathrm{P}\) the complex Lie subgroup of \(\mathrm{G}_c\) with Lie algebra \(\mathfrak{p}\) . Show that the canonical map \(\mathrm{G / H} \to \mathrm{G}_c / \mathrm{P}\) is an isomorphism of complex-analytic manifolds.*

\begin{enumerate}
\def\labelenumi{\arabic{enumi})}
\setcounter{enumi}{8}
\tightlist
\item
  Let H be a connected closed subgroup of G, of maximum rank, distinct from G and maximal for these properties. Denote by Z the quotient group \(\mathrm{C}(\mathrm{H})/\mathrm{C}(\mathrm{G})\) .
\end{enumerate}

\begin{enumerate}
\def\labelenumi{\alph{enumi})}
\item
  We have \(\dim Z \leq 1\) ; if \(\dim Z = 0\) , then Z is of order 2,3 or 5 (reduce to the case in which G is almost simple, with trivial centre; apply Chap. VI, §4, Exerc. 4).
\item
  Assume that G is almost simple and that Z is of order 3 or 5; put \(z = \text{Card}(Z)\) and \(\pi = \text{Card}(\pi_1(H))/\text{Card}(\pi_1(G))\) . Show that we must be in one of the following seven cases:
\end{enumerate}

\begin{enumerate}
\def\labelenumi{(\roman{enumi})}
\item
  G is of type \(\mathrm{G}_2\) , H is of type \(\mathrm{A}_2\) , and \(z = 3, \pi = 1\) ;
\item
  G is of type \(F_{4}\) , H is of type \(A_{2} \times A_{2}\) , and \(z = \pi = 3\) ;
\item
  G is of type \(E_{6}\) , H is of type \(A_{2} \times A_{2} \times A_{2}\) , and \(z = \pi = 3\) ;
\item
  G is of type \(E_{7}\) , H is of type \(A_{2} \times A_{5}\) , and \(z = \pi = 3\) ;
\item
  G is of type \(E_{8}\) , H is of type \(A_{8}\) , and \(z = \pi = 3\) ;
\item
  G is of type \(E_{8}\) , H is of type \(A_{2} \times E_{6}\) , and \(z = \pi = 3\) ;
\item
  G is of type \(E_{8}\) , H is of type \(A_{4} \times A_{4}\) , and \(z = \pi = 5\) .
\end{enumerate}

(Use loc. cit. and the plates of Chap. VI; to calculate \(\pi\) , remark that if \(f\) and \(f'\) denote the connection indices of G and H, respectively, then \(z\pi f = f'\) .)

\begin{enumerate}
\def\labelenumi{\alph{enumi})}
\setcounter{enumi}{2}
\tightlist
\item
  In each of the preceding cases, determine the group H.
\end{enumerate}

\begin{enumerate}
\def\labelenumi{\arabic{enumi})}
\setcounter{enumi}{9}
\tightlist
\item
  We retain the notations of the preceding exercise.
\end{enumerate}

\begin{enumerate}
\def\labelenumi{\alph{enumi})}
\item
  Assume that \(\dim Z = 1\) . Then there are exactly two complex structures on G/H invariant under G; there exists an automorphism of G that leaves H stable and interchanges these two structures (use Exerc. 8).
\item
  Determine the complex structures on \(\mathbf{P}_n(\mathbf{C})\) invariant under \(\mathbf{SU}(n + 1,\mathbf{C})\) .
\item
  Assume that \(\dim Z = 0\) and \(\operatorname{Card}(Z) \neq 2\) . Show that there exist almost complex structures on G/H invariant under G (if z is an element of C(H) that is not central in G, Int z induces an automorphism of G/H of odd order (Exerc. 9); use Exerc. 8 b)). Show that there is no complex structure on G/H invariant under G (use Exerc. 8 d)).
\end{enumerate}

\(d\) ) There is no complex structure on \(\mathbf{S}_6\) invariant under \(\mathbf{SO}(7,\mathbf{R})\) (use Exerc. 7 of §3).

\begin{enumerate}
\def\labelenumi{\arabic{enumi})}
\setcounter{enumi}{10}
\tightlist
\item
  We retain the notations of Exerc. 9 and 10.
\end{enumerate}

\begin{enumerate}
\def\labelenumi{\alph{enumi})}
\item
  The homogeneous space G/H is symmetric (§1, Exerc. 8) if and only if Z is of order 2 or of dimension 1. The symmetric space G/H is then irreducible (loc. cit.).
\item
  Assume that \(\dim Z = 1\) ; denote by \(X\) the complex-analytic manifold \(G / H\) . Show that we are then in one of the following situations:
\end{enumerate}

\begin{enumerate}
\def\labelenumi{(\roman{enumi})}
\item
  \(\mathbf{G}\) is of type \(\mathrm{A}_l\) and \(\mathrm{D(H)}\) of type \(\mathrm{A}_{p - 1}\times \mathrm{A}_{l - p}\) ; \(\mathbf{X}\) is isomorphic to the grassmannian \(\mathbf{G}_p(\mathbf{C}^{l + 1})\) .
\item
  G is of type \(B_{l}\) and D(H) of type \(B_{l-1}\) ; X is isomorphic to the submanifold of \(\mathbf{P}_{2l}(\mathbf{C})\) defined by the vanishing of a separating quadratic form (smooth projective quadric).
\end{enumerate}

(ii') G is of type \(D_{l}\) and \(D(H)\) of type \(D_{l-1}\) ; X is isomorphic to a smooth projective quadric in \(\mathbf{P}_{2l-1}(\mathbf{C})\) .

\begin{enumerate}
\def\labelenumi{(\roman{enumi})}
\setcounter{enumi}{2}
\item
  G is of type \(C_{l}\) and D(H) of type \(A_{l-1}\) ; X is isomorphic to the submanifold of \(\mathbf{G}_{l}(\mathbf{C}^{2l})\) consisting of the maximal isotropic subspaces for the usual alternating bilinear form on \(C^{2l}\) .
\item
  G is of type \(D_{l}\) and \(\mathrm{D}(\mathrm{H})\) of type \(A_{l-1}\) ; X is isomorphic to the submanifold of \(\mathbf{G}_{l}(\mathbf{C}^{2l})\) consisting of the maximal isotropic subspaces for the usual symmetric bilinear form on \(C^{2l}\) .
\item
  G is of type \(E_{6}\) and D(H) of type \(D_{5}\) .
\item
  G is of type \(E_{7}\) and D(H) of type \(E_{6}\) .
\end{enumerate}

\begin{enumerate}
\def\labelenumi{\alph{enumi})}
\setcounter{enumi}{2}
\tightlist
\item
  Assume that \(\operatorname{Card}(Z)=2\) ; give the list of possible situations. If G is of type \(A_{l}\) , \(B_{l}\) , \(C_{l}\) or \(D_{l}\) , show that the real manifold G/H is isomorphic to a grassmannian manifold \(\mathrm{G}_{p}(\mathrm{K}^{q})\) , with K=R, C or H.
\end{enumerate}

¶ 12) Assume that G is simply-connected; for all \(\alpha \in \mathrm{R}(\mathrm{G},\mathrm{T})\) , put \(t_{\alpha} = \exp (\frac{1}{2} K_{\alpha})\) . Put \(\mathrm{N} = \mathrm{N}_{\mathrm{G}}(\mathrm{T})\) , and denote by \(\varphi : \mathrm{N} \to \mathrm{W}\) the canonical map. Let B be a basis of \(\mathrm{R}(\mathrm{G},\mathrm{T})\) ; choose an element \(n_{\alpha}\) of \((\mathrm{N} \cap \mathrm{S}_{\alpha}) - (\mathrm{T} \cap \mathrm{S}_{\alpha})\) for all \(\alpha \in \mathrm{B}\) .

\begin{enumerate}
\def\labelenumi{\alph{enumi})}
\item
  Show that \(\varphi(n_{\alpha}) = s_{\alpha}\) and \(n_{\alpha}^{2} = t_{\alpha}\) , so that \(n_{\alpha}^{4} = 1\) .
\item
  Let \(\alpha\) and \(\beta\) be two distinct elements of B, and let \(m_{\alpha \beta}\) be the order of \(s_{\alpha}s_{\beta}\) in W. Prove that
\end{enumerate}

\[
\begin{array}{r l} {n _ {\alpha} n _ {\beta} = n _ {\beta} n _ {\alpha}} & {\mathrm{if} m _ {\alpha \beta} = 2} \\ {n _ {\alpha} n _ {\beta} n _ {\alpha} = n _ {\beta} n _ {\alpha} n _ {\beta}} & {\mathrm{if} m _ {\alpha \beta} = 3} \\ {(n _ {\alpha} n _ {\beta}) ^ {2} = (n _ {\beta} n _ {\alpha}) ^ {2}} & {\mathrm{if} m _ {\alpha \beta} = 4} \\ {(n _ {\alpha} n _ {\beta}) ^ {3} = (n _ {\beta} n _ {\alpha}) ^ {3}} & {\mathrm{if} m _ {\alpha \beta} = 6} \end{array}
\]

(if for example \(m_{\alpha\beta}=3\) , then \((s_{\alpha}s_{\beta})s_{\alpha}(s_{\alpha}s_{\beta})^{-1}=s_{\beta}\) and \(s_{\alpha}s_{\beta}(\alpha)=\beta\) ; show that \(n_{\alpha}n_{\beta}n_{\alpha}n_{\beta}^{-1}n_{\alpha}^{-1}n_{\beta}^{-1}\) belongs to \(S_{\beta}\) , and conclude by remarking that \(S_{\alpha}\cap S_{\beta}\cap T=\{e\}\) ).

\begin{enumerate}
\def\labelenumi{\alph{enumi})}
\setcounter{enumi}{2}
\item
  Deduce from b) that there exists a unique section \(\nu: W \to N\) of \(\varphi\) such that \(\nu(s_{\alpha}) = n_{\alpha}\) and \(\nu(ww') = \nu(w)\nu(w')\) if \(l(ww') = l(w) + l(w')\) (where \(l(w)\) denotes the length of w with respect to the generating system \((s_{\alpha})_{\alpha \in \mathbb{B}}\) ; cf.~Chap. VI, §1, no. 5, Prop. 5). Put \(n_w = \nu(w)\) .
\item
  Let \(W^{*}\) be the subgroup of N generated by the \(n_{\alpha}\) ; show that \(W^{*} \cap T\) is the subgroup \(T_{2}\) of T consisting of the elements of order \(\leq 2\) and that W can be identified with \(W^{*}/T_{2}\) .
\item
  Let \(w \in \mathrm{W}\) be such that \(w^2 = 1\) ; show that \(n_w^2 = \prod_{\alpha \in \mathrm{R}_w} t_\alpha\) , where \(\mathrm{R}_w\) is the set of positive roots \(\alpha\) such that \(w(\alpha) < 0\) (write \(w = s_{\alpha_1} \ldots s_{\alpha_r}\) , with \(r = l(w)\) and \(\alpha_i \in \mathrm{B}\) ; apply c) and Chap. VI, §1, no. 6, Cor. 2 of Prop. 17).
\end{enumerate}

\(f\) ) Assume that \(\mathrm{G}\) is almost simple. Let \(c\) be a Coxeter transformation in \(\mathrm{W}\) and \(h\) the Coxeter number of \(\mathrm{W}\) (Chap. VI, §1, no. 11). Show that \(n_c^h = \prod_{\alpha \in \mathbb{R}_+}t_\alpha\) (use \(c\) ), \(e\) ) and Exerc. 2 of Chap. V, §6).

\begin{enumerate}
\def\labelenumi{\arabic{enumi})}
\setcounter{enumi}{12}
\tightlist
\item
  \begin{enumerate}
  \def\labelenumii{\alph{enumii})}
  \tightlist
  \item
    Choose a basis B of R(G, T), and denote by R+ the set of positive roots of R(G, T). Prove that \(z_{G} = \prod_{\alpha \in R_{+}} \exp(\frac{1}{2} K_{\alpha})\) is an element of C(G), independent of the choice of T and of B; we have \(z_{G}^{2} = e\) .
  \end{enumerate}
\end{enumerate}

\begin{enumerate}
\def\labelenumi{\alph{enumi})}
\setcounter{enumi}{1}
\item
  Let H be another connected compact Lie group. Then \(z_{\mathrm{G}\times\mathrm{H}} = (z_{\mathrm{G}}, z_{\mathrm{H}})\) ; prove that, if \(f : G \to H\) is a surjective morphism of Lie groups, then \(f(z_{\mathrm{G}}) = z_{\mathrm{H}}\) .
\item
  Put \(\mathrm{R} = \mathrm{R}(\mathrm{G},\mathrm{T})\) ; assume that \(\mathbf{G}\) is simply-connected, so that \(\mathrm{X(T)}\) can be identified with \(\mathrm{P(R)}\) . Show that the kernel of the homomorphism \(\chi \mapsto \chi (z_{\mathrm{G}})\)
\end{enumerate}

from X(T) to \(\{1,-1\}\) is the subgroup \(P'(R)\) defined in Exerc. 8 of Chap. VI, §1.

\begin{enumerate}
\def\labelenumi{\alph{enumi})}
\setcounter{enumi}{3}
\tightlist
\item
  If \(\mathrm{G} = \mathbf{S}\mathbf{U}(n,\mathbf{C})\) , then \(z_{\mathrm{G}} = (-1)^{n + 1}I_n\) ;
\end{enumerate}

if \(\mathrm{G} = \mathbf{S}\mathbf{U}(n,\mathbf{H})\) , then \(z_{\mathrm{G}} = -I_n\) ;

if \(\mathrm{G} = \mathbf{Spin}(n, \mathbf{R})\) with \(n \equiv 3, 4, 5\) or 6 (mod 8), then \(z_{G} = -1\) ;

if \(\mathrm{G} = \mathbf{Spin}(n, \mathbf{R})\) with \(n \equiv 0, 1, 2, 7 \pmod{8}\) , then \(z_{G} = 1\) ;

if G is of type \(E_{6}\) , \(E_{8}\) , \(F_{4}\) or \(G_{2}\) , then \(z_{G} = e\) ;

if G is simply-connected of type \(E_{7}\) , then \(z_{G}\) is the unique non-identity element of \(\mathrm{C}(\mathrm{G})\) .

(Use Chap. VI, §4, Exerc. 5.)

\begin{enumerate}
\def\labelenumi{\arabic{enumi})}
\setcounter{enumi}{13}
\tightlist
\item
  Assume that the group G is almost simple; denote the Coxeter number of R(G,T) by h (Chap. VI, §1, no. 11). An element g of G is said to be a Coxeter element if there exists a maximal torus S of G such that g belongs to \(\mathrm{N}_{\mathrm{G}}(\mathrm{S})\) and its class in \(\mathrm{W}_{\mathrm{G}}(\mathrm{S})\) is a Coxeter transformation (loc. cit.).
\end{enumerate}

\begin{enumerate}
\def\labelenumi{\alph{enumi})}
\item
  Show that any two Coxeter elements are conjugate (argue as in the proof of the Cor. of Prop. 10, no. 5).
\item
  Every Coxeter element g is regular and satisfies \(g^{h} = z_{G}\) , where \(z_{G}\) is the element of C(G) defined in Exerc. 13; in particular, g is of order h or 2h according as \(z_{G}\) is or is not equal to e (use Exerc. 12 f)).
\item
  An element \(g \in G\) is a Coxeter element if and only if the automorphism \(Ad g \otimes 1_{C}\) of \(g_{C}\) satisfies the equivalent conditions of Chap. VIII, §5, Exerc. 5 f).
\end{enumerate}

\(d\) ) Show that every regular element \(g\) of \(\mathbf{G}\) such that \(g^{h}\in \mathrm{C}(\mathrm{G})\) is a Coxeter element; for \(p < h\) , there is no regular element \(k\) such that \(k^p\in \mathrm{C}(\mathrm{G})\) .

\begin{enumerate}
\def\labelenumi{\arabic{enumi})}
\setcounter{enumi}{14}
\tightlist
\item
  Let H be a connected closed subgroup of G. Then H is said to be clean if it is not contained in any connected closed subgroup of maximal rank distinct from G.
\end{enumerate}

\begin{enumerate}
\def\labelenumi{\alph{enumi})}
\item
  Show that H is clean if and only if its centralizer in G is equal to C(G). In particular, if H is clean then C(H) = C(G) ∩ H.
\item
  Assume from now on that H is clean. Show that, for every maximal torus S of H, we have \(\mathrm{C}(\mathrm{H}) = \mathrm{S} \cap \mathrm{C}(\mathrm{G})\) .
\item
  Let \(H'\) be a connected closed subgroup of G containing H. Prove that \(rkH < rkH'\) , and deduce that H is clean in \(H'\) .
\item
  Let K be a connected closed subgroup of H, clean in H and containing a regular element of G. Show that K is clean in G.
\end{enumerate}

\begin{enumerate}
\def\labelenumi{\arabic{enumi})}
\setcounter{enumi}{15}
\tightlist
\item
  Let H be a connected closed subgroup of G such that \(T \cap H\) is a maximal torus S of H. Let \(\lambda \in \mathrm{R}(\mathrm{H}, \mathrm{S})\) ; denote by \(\mathrm{R}(\lambda)\) the set of roots of \(\mathrm{R}(\mathrm{G}, \mathrm{T})\) whose restriction to S is equal to \(\lambda\) .
\end{enumerate}

\begin{enumerate}
\def\labelenumi{\alph{enumi})}
\item
  Show that \(\mathrm{R}(\lambda)\) is not empty.
\item
  Let \(w \in \mathrm{W}_{\mathrm{H}}(\mathrm{S})\) ; show that there exists an element \(\bar{w}\) of \(\mathrm{W}_{\mathrm{G}}(\mathrm{T})\) such that \(\mathrm{R}(w\lambda) = \bar{w}\mathrm{R}(\lambda)\) (use Exerc. 4 of §2).
\item
  Let \(\mathrm{P}(\lambda)\) be the intersection of \(\mathrm{R}(\mathrm{G},\mathrm{T})\) with the subgroup of \(\mathrm{X}(\mathrm{T})\) generated by \(\mathrm{R}(\lambda)\) . Show that there exists a closed subgroup \(G_{\lambda}\) of G containing T such that \(\mathrm{P}(\lambda)=\mathrm{R}(\mathrm{G}_{\lambda},\mathrm{T})\) . Deduce that the reflection \(s_{\lambda}\in\mathrm{W}_{\mathrm{H}}(\mathrm{S})\) is the restriction to S of a product of reflections \(s_{\alpha}\) with \(\alpha\in\mathrm{P}(\lambda)\) .
\item
  Show that the nodal vector \(K_{\lambda} \in \mathrm{L}(\mathrm{S})\) associated to \(\lambda\) is a linear combination with integer coefficients of the \(K_{\alpha}\) for \(\alpha \in \mathrm{R}(\lambda)\) .
\item
  Let \(B_{H}\) be a basis of \(R(H,S)\) . Show that \(R(\lambda)\) is contained in the subgroup of \(X(T)\) generated by the union of the \(R(\mu)\) for \(\mu \in B_{H}\) (prove that the set of \(\lambda \in R(H,S)\) having the stated property is stable under \(s_{\mu}\) for all \(\mu \in B_{H}\) , and use c)).
\end{enumerate}

¶ 17) We retain the notations of the preceding exercise; assume further that the subgroup H is clean (Exerc. 15).

\begin{enumerate}
\def\labelenumi{\alph{enumi})}
\item
  Show that \(\mathrm{R}(\mathrm{G},\mathrm{T})\) is contained in the subgroup of \(\mathrm{X}(\mathrm{T})\) generated by the union of the \(\mathrm{R}(\lambda)\) for \(\lambda\in B_{H}\) .
\item
  Let \(\Delta\) be the identity component of the subgroup of S consisting of the \(s \in S\) such that \(\lambda(s) = \mu(s)\) for all \(\lambda, \mu\) in \(B_{H}\) . Show that there exists a basis \(\{\alpha_{1}, \ldots, \alpha_{l}\}\) of R(G,T) and an integer k, with \(0 \leq k \leq l - 1\) , such that \(\Delta\) is the identity component of the set of \(t \in T\) satisfying
\end{enumerate}

\[
\alpha_ {1} (t) = \dots = \alpha_ {k} (t) = 1, \alpha_ {k + 1} (t) = \dots = \alpha_ {l} (t).
\]

(Let \(x\) be the element of \(\mathrm{L}(\mathrm{S})\) such that \(\delta(\lambda)(x) = 2\pi i\) for all \(\lambda \in \mathrm{B}_{\mathrm{H}}\) ; deduce from \(a\) ) that \(\exp x \in \mathrm{C}(\mathrm{G})\) . Choose \(\{\alpha_1, \ldots, \alpha_l\}\) so that \(i\delta(\alpha_j)(x)\) is zero for \(1 \leq j \leq k\) and \(< 0\) for \(k + 1 \leq j \leq l\) ; then show that for all \(\lambda \in \mathrm{B}_{\mathrm{H}}\) every root in \(\mathrm{R}(\lambda)\) can be written

\[
\alpha_ {j} + n _ {1} \alpha_ {1} + \dots + n _ {k} \alpha_ {k},
\]

with \(j \geq k + 1\) and \(n_i \in \mathbf{N}\) . Conclude by using \(a)\) .

\begin{enumerate}
\def\labelenumi{\alph{enumi})}
\setcounter{enumi}{2}
\tightlist
\item
  Show that the integer k does not depend on the choice of the tori S, T or of the bases \(B_{H}, \{\alpha_{1}, \ldots, \alpha_{l}\}\) ; it is equal to the rank of the derived group of \(Z_{\mathrm{G}}(\Delta)\) .
\end{enumerate}

\begin{enumerate}
\def\labelenumi{\arabic{enumi})}
\setcounter{enumi}{17}
\tightlist
\item
  We retain the notations of Exerc. 16 and 17. The subgroup H is said to be principal if it is clean and if the sub-torus \(\Delta\) contains a regular element (in other words, if k = 0).
\end{enumerate}

\begin{enumerate}
\def\labelenumi{\alph{enumi})}
\item
  Let \(\mathrm{K}\) be a connected closed subgroup of \(\mathrm{H}\) . Show that \(\mathrm{K}\) is a principal subgroup of \(\mathrm{G}\) if and only if \(\mathrm{K}\) is principal in \(\mathrm{H}\) and \(\mathrm{H}\) is principal in \(\mathrm{G}\) (use Exerc. 15 c) and d)).
\item
  Assume from now on that H is principal. Show that the union of the \(\mathrm{R}(\mu)\) , for \(\mu \in B_{H}\) , is a basis of \(\mathrm{R}(\mathrm{G}, \mathrm{T})\) .
\item
  Let \(\alpha \in \mathrm{R}(\mathrm{G},\mathrm{T})\) . Prove that the restriction of \(\alpha\) to S is a non-zero integer multiple of a root \(\lambda \in \mathrm{R}(\mathrm{H},\mathrm{S})\) ; the root \(\alpha\) is a sum of elements of \(\mathrm{R}(\lambda)\) (show that the intersection of \(\mathrm{L}(\mathrm{S})\) with the hyperplane \(\delta (\alpha) = 0\) is a wall of \(\mathrm{L}(\mathrm{S}))\) ).
\item
  Let \(\lambda\in\mathrm{R}(\mathrm{H},\mathrm{S})\) ; show that the nodal vector \(K_{\lambda}\) is a linear combination with non-negative integer coefficients of the \(K_{\alpha}\) for \(\alpha\in\mathrm{R}(\lambda)\) (cf.~Exerc. 16 d) and Chap. V, §3, no. 5, Lemma 6).
\end{enumerate}

\begin{enumerate}
\def\labelenumi{\arabic{enumi})}
\setcounter{enumi}{18}
\tightlist
\item
  We retain the notations of Exerc. 16 to 18; assume that the subgroup H is principal. Give \(\mathrm{X(T)\otimes\mathbf{R}}\) (resp. \(\mathrm{X(S)\otimes\mathbf{R}}\) ) a scalar product invariant under \(\mathrm{W_{G}(T)}\) (resp. under \(\mathrm{W_{H}(S)}\) ). Let \(\lambda,\mu\) be two roots in \(B_{H}\) .
\end{enumerate}

\begin{enumerate}
\def\labelenumi{\alph{enumi})}
\item
  Show that, if \(\lambda\) and \(\mu\) are orthogonal, the sets \(\mathrm{R}(\lambda)\) and \(\mathrm{R}(\mu)\) are orthogonal. Deduce that if \(\mathrm{G}\) is almost simple then so is \(\mathrm{H}\) .
\item
  Assume from now on that \(n(\lambda, \mu) = -1\) . Show that there exists a surjective map \(u: \mathrm{R}(\lambda) \to \mathrm{R}(\mu)\) such that, for \(\alpha \in \mathrm{R}(\lambda)\) , \(u(\alpha)\) is the unique root of \(\mathrm{R}(\mu)\) linked (that is, not orthogonal) to \(\alpha\) ; we have \(K_{\mu} = \sum_{\beta \in \mathrm{R}(\mu)} K_{\beta}\) , and the roots in \(\mathrm{R}(\mu)\) are pairwise orthogonal (write \(K_{\mu}\) and \(K_{\lambda}\) as linear combinations of the \(K_{\alpha}\) , then identify the coefficients).
\item
  If \(n(\mu, \lambda) = -1\) , the map \(u\) is bijective, and the only pairs of linked roots in \(\mathrm{R}(\lambda) \cup \mathrm{R}(\mu)\) are the pairs \((\alpha, u(\alpha))\) for \(\alpha \in \mathrm{R}(\lambda)\) .
\item
  Assume that \(n(\mu,\lambda)=-2\) ; let \(\beta\in\mathrm{R}(\mu)\) . Show that \(u^{-1}(\beta)\) contains either one or two elements; if \(u^{-1}(\beta)=\{\alpha_{1},\alpha_{2}\}\) , the root \(\alpha_{i}\) (i=1,2) is orthogonal to the other roots in \(\mathrm{R}(\lambda)\) , and has the same length as \(\beta\) ; if \(u^{-1}(\beta)=\{\alpha\}\) and \(\|\alpha\|=\|\beta\|\) , \(\alpha\) is linked to a unique root \(\alpha^{\prime}\in\mathrm{R}(\lambda)\) , of the same length as \(\alpha\) , and \(\{\alpha,\alpha^{\prime}\}\) is orthogonal to the remainder of \(\mathrm{R}(\lambda)\) ; if \(u^{-1}(\beta)=\{\alpha\}\) and \(\|\alpha\|\neq\|\beta\|\) , \(\alpha\) is orthogonal to the other roots in \(\mathrm{R}(\lambda)\) .
\item
  Study similarly the case \(n(\mu, \lambda) = -3\) .
\end{enumerate}

\begin{enumerate}
\def\labelenumi{\arabic{enumi})}
\setcounter{enumi}{19}
\tightlist
\item
  Assume that the group G is almost simple; let H be a principal subgroup of G, of rank \(\geq 2\) .
\end{enumerate}

\begin{enumerate}
\def\labelenumi{\alph{enumi})}
\item
  Show that H is semi-simple of type \(B_{h}, C_{h}, F_{4}\) or \(G_{2}\) (use Exerc. 19).
\item
  Assume that H is of rank \(\geq 3\) . Show that G is of type \(A_{l}, D_{l}, E_{6}, E_{7}\) or \(E_{8}\) (consider the terminal vertices of the Dynkin graph of R(G, T), and apply Exerc. 19).
\item
  Show that, if \(rkH \geq 3\) , we are in one of the following situations:
\end{enumerate}

G is of type \(A_{2l}\) \((l \geq 3)\) and H of type \(B_{l}\) ;

G is of type \(\mathrm{A}_{2l - 1}\) \((l\geq 3)\) and H of type \(\mathbf{C}_l\)

G is of type \(D_{l}\) \((l \geq 4)\) and H of type \(B_{l-1}\) ;

G is of type \(\mathrm{E}_6\) and H of type \(\mathrm{F}_4\) .

\(d\) ) Show that, if H is of type \(\mathrm{B}_2\) , then G is of type \(\mathrm{A}_3\) or \(\mathrm{A}_4\) .

\begin{enumerate}
\def\labelenumi{\alph{enumi})}
\setcounter{enumi}{4}
\tightlist
\item
  Show that, if H is of type \(G_{2}\) , then G is of type \(B_{3}, D_{4}\) or \(A_{6}\) .
\end{enumerate}

(For a more explicit description of these situations, see §5, Exerc. 5.)

\(f)\) Let \(\mathrm{K}\) be a connected closed subgroup of \(\mathrm{G}\) containing \(\mathrm{H}\) ; show that either \(\mathrm{K} = \mathrm{G}\) or \(\mathrm{K} = \mathrm{H}\) , or \(\mathrm{H}\) is of type \(\mathrm{G}_2\) , \(\mathrm{K}\) of type \(\mathrm{B}_3\) and \(\mathrm{G}\) of type \(\mathrm{D}_4\) or \(\mathrm{A}_6\) .

\begin{enumerate}
\def\labelenumi{\arabic{enumi})}
\setcounter{enumi}{20}
\tightlist
\item
  \begin{enumerate}
  \def\labelenumii{\alph{enumii})}
  \tightlist
  \item
    Let H be a connected closed subgroup of G of rank 1. Show that the following conditions are equivalent:
  \end{enumerate}
\end{enumerate}

\begin{enumerate}
\def\labelenumi{(\roman{enumi})}
\item
  H is principal;
\item
  H is clean and contains a regular element of G;
\item
  there exists a principal \(\mathfrak{sl}_2\) -triplet \((x,h,y)\) in \(\mathfrak{g}_{\mathbf{C}}\) (Chap. VIII, §11, no. 4) such that
\end{enumerate}

\[
\mathrm{L} (\mathrm{H}) _ {(\mathbf {C})} = \mathbf {C} x + \mathbf {C} h + \mathbf {C} y.
\]

\begin{enumerate}
\def\labelenumi{\alph{enumi})}
\setcounter{enumi}{1}
\item
  Show that G contains a principal subgroup of rank 1, and that any two such subgroups are conjugate (with the notations of Exerc. 17, remark that \(S = \Delta\) ).
\item
  Show that a connected closed subgroup of G is principal if and only if it contains a principal subgroup of G of rank 1 (use Exerc. 18 a)).
\end{enumerate}

\begin{enumerate}
\def\labelenumi{\arabic{enumi})}
\setcounter{enumi}{21}
\item
  Let H be a principal subgroup of G of rank 1; let \(\Gamma\) be the subgroup of \(\operatorname{Aut}(G)\) consisting of the automorphisms u such that \(u(\mathrm{H}) = \mathrm{H}\) . Show that \(\operatorname{Aut}(G) = \Gamma.\operatorname{Int}(G)\) .
\end{enumerate}

